\section{Limitations}
\label{sec:limitations}

The conclusions should be read within the scope of the available automatic metrics, internal baselines, and limited perceptual evidence.
The ESD SSST and DSST results provide the clearest evidence for an alpha-conditioned emotion--speaker frontier, while CREMA-D from scratch is better interpreted as a stress test than as a stable cross-dataset confirmation.
The current CREMA-D from-scratch results also show label-specific instability, especially for the native-ranker epoch-30 condition, so they should not be used to claim robust CREMA-D emotion control.

The speaker metrics are based on speaker embeddings and cosine similarity, which are useful but not identical to human judgments of speaker identity.
Likewise, the emotion-side automatic metrics depend on a frozen emotion model and should be paired with listening tests before making strong perceptual claims.
The available listening study supports the ESD SSST frontier direction, but it does not yet cover DSST, CREMA-D, or all ablation branches.

The model-comparison significance tests are paired over samples and alpha values, not over multiple independent training seeds.
They support differences between evaluated outputs under the tested checkpoints, but they do not prove that a specific loss term is causally responsible across random seeds or datasets.
Future work should add seed-level replication, perceptual tests for the DSST and CREMA-D conditions, and external EVC systems beyond the internal StarGAN-style and CycleGAN baselines.
